\documentclass[10pt,a4paper]{amsart}
\usepackage{amsmath,amssymb,amsthm,graphicx}
% Retained minimal preamble for the Stacks Project chapter more-morphisms.tex (More on Morphisms).
% Verbatim excerpt of the chapter preamble: Xy-pic setup (native preamble lines 14-18) and the native
% theorem declarations with the mathematical macros (native preamble lines 154-217). Copied mechanically;
% nothing added, removed or redefined. The wrapper class and the amsmath/amssymb/amsthm/graphicx base
% packages are supplied by the assembler, so the class/externaldocument part of the native preamble is omitted.
% For commutative diagrams we use Xy-pic
\usepackage[all]{xy}

% We use 2cell for 2-commutative diagrams.
\xyoption{2cell}
% Theorem environments.
%
\theoremstyle{plain}
\newtheorem{theorem}[subsection]{Theorem}
\newtheorem{proposition}[subsection]{Proposition}
\newtheorem{lemma}[subsection]{Lemma}

\theoremstyle{definition}
\newtheorem{definition}[subsection]{Definition}
\newtheorem{example}[subsection]{Example}
\newtheorem{exercise}[subsection]{Exercise}
\newtheorem{situation}[subsection]{Situation}

\theoremstyle{remark}
\newtheorem{remark}[subsection]{Remark}
\newtheorem{remarks}[subsection]{Remarks}

\numberwithin{equation}{subsection}

% Macros
%
\def\lim{\mathop{\mathrm{lim}}\nolimits}
\def\colim{\mathop{\mathrm{colim}}\nolimits}
\def\Spec{\mathop{\mathrm{Spec}}}
\def\Hom{\mathop{\mathrm{Hom}}\nolimits}
\def\Ext{\mathop{\mathrm{Ext}}\nolimits}
\def\SheafHom{\mathop{\mathcal{H}\!\mathit{om}}\nolimits}
\def\SheafExt{\mathop{\mathcal{E}\!\mathit{xt}}\nolimits}
\def\Sch{\mathit{Sch}}
\def\Mor{\mathop{\mathrm{Mor}}\nolimits}
\def\Ob{\mathop{\mathrm{Ob}}\nolimits}
\def\Sh{\mathop{\mathit{Sh}}\nolimits}
\def\NL{\mathop{N\!L}\nolimits}
\def\CH{\mathop{\mathrm{CH}}\nolimits}
\def\proetale{{pro\text{-}\acute{e}tale}}
\def\etale{{\acute{e}tale}}
\def\QCoh{\mathit{QCoh}}
\def\Ker{\mathop{\mathrm{Ker}}}
\def\Im{\mathop{\mathrm{Im}}}
\def\Coker{\mathop{\mathrm{Coker}}}
\def\Coim{\mathop{\mathrm{Coim}}}

% Boxtimes
%
\DeclareMathSymbol{\boxtimes}{\mathbin}{AMSa}{"02}

%
% Macros for moduli stacks/spaces
%
\def\QCohstack{\mathcal{QC}\!\mathit{oh}}
\def\Cohstack{\mathcal{C}\!\mathit{oh}}
\def\Spacesstack{\mathcal{S}\!\mathit{paces}}
\def\Quotfunctor{\mathrm{Quot}}
\def\Hilbfunctor{\mathrm{Hilb}}
\def\Curvesstack{\mathcal{C}\!\mathit{urves}}
\def\Polarizedstack{\mathcal{P}\!\mathit{olarized}}
\def\Complexesstack{\mathcal{C}\!\mathit{omplexes}}
% \Pic is the operator that assigns to X its picard group, usage \Pic(X)
% \Picardstack_{X/B} denotes the Picard stack of X over B
% \Picardfunctor_{X/B} denotes the Picard functor of X over B
\def\Pic{\mathop{\mathrm{Pic}}\nolimits}
\def\Picardstack{\mathcal{P}\!\mathit{ic}}
\def\Picardfunctor{\mathrm{Pic}}
\def\Deformationcategory{\mathcal{D}\!\mathit{ef}}
\usepackage{cleveref}
\crefname{originalreference}{original reference}{original references}
\Crefname{originalreference}{Original reference}{Original references}
\makeatletter
\newlabel{morphisms-lemma-quasi-compact-scheme-theoretic-image}{{01R8}{source}{}{sourcepointer-1}{}}
\newlabel{morphisms-lemma-quasi-compact-scheme-theoretic-image@cref}{{[originalreference][0][]01R8}{[1][0][]source}{}{}{}}
\makeatother
\crefname{originalreference}{original reference}{original references}
\Crefname{originalreference}{Original reference}{Original references}
\makeatletter
\newlabel{schemes-lemma-morphism-into-affine}{{01I1}{source}{}{sourcepointer-2}{}}
\newlabel{schemes-lemma-morphism-into-affine@cref}{{[originalreference][0][]01I1}{[1][0][]source}{}{}{}}
\makeatother
\title{The canonical morphism to the spectrum of global functions is dominant over quasi-compact schemes}
\author{The Stacks Project Authors}
\date{stacks-project:more-morphisms.tex:local-snapshot-20261011; GNU FDL 1.2 or later}
\begin{document}
\maketitle
\noindent\textit{Editorial note.} Counted span is the author's own proof environment, source lines 23577-23596, immediately following the statement at 23565-23575.\par\medskip
\noindent\textit{Editorial note.} Two cross-chapter pointers (Morphisms, Schemes) are recorded in assembly.external\_references.\par\medskip
\noindent\textit{Editorial note.} The statement's three assertions (vanishing of the largest quasi-coherent submodule, injectivity, dominance) are all addressed inside the span.\par\medskip
\noindent\textit{Allowed standard prerequisites (editorial).}\begin{itemize}
\item The canonical morphism into the spectrum of global sections. Schemes, lemma schemes-lemma-morphism-into-affine, cited in the statement. Fixes the morphism f : X -> S = Spec(Gamma(X, O\_X)) whose injectivity is studied.
\item Scheme-theoretic image of a quasi-compact morphism. Morphisms, lemma morphisms-lemma-quasi-compact-scheme-theoretic-image, cited in the proof. Identifies the closed subscheme cutting out the kernel, which is then shown to be empty.
\item Quasi-coherent modules on an affine scheme. used to pass from the kernel sheaf to an ideal of the ring of global functions. Turns the kernel of f\textasciicircum{}sharp into an ideal, so that vanishing is visible on the quasi-compact pieces.
\end{itemize}\medskip
\ifdefined\hypertarget\hypertarget{sourcepointer-1}{}\fi
\noindent\textit{Source reference.} Original reference 01R8: focused-seeds/geometry/S9-morph-aff-injective-on-sheaves/ref-tags.txt, source lines 4--4. Navigation pointer only: the counted proof cites this statement by \textbackslash{}ref as a standing input; the deduction the counted proof has to make is carried out inside the counted span, and no part of that deduction is delegated to this pointer.\par\medskip
\ifdefined\hypertarget\hypertarget{sourcepointer-2}{}\fi
\noindent\textit{Source reference.} Original reference 01I1: focused-seeds/geometry/S9-morph-aff-injective-on-sheaves/ref-tags.txt, source lines 5--5. Navigation pointer only: the counted proof cites this statement by \textbackslash{}ref as a standing input; the deduction the counted proof has to make is carried out inside the counted span, and no part of that deduction is delegated to this pointer.\par\medskip
\begin{lemma}
\label{lemma-aff-is-injective-on-sheaves}
Let $X$ be a scheme. Let $f : X \to S = \Spec(\Gamma(X, \mathcal{O}_X))$
be the canonical morphism of
Schemes, Lemma \ref{schemes-lemma-morphism-into-affine}.
The largest quasi-coherent $\mathcal{O}_S$-module contained
in the kernel of $f^\sharp : \mathcal{O}_S \to f_*\mathcal{O}_X$
is zero. If $X$ is quasi-compact, then $f^\sharp$ is injective.
In particular, if $X$ is quasi-compact, then $f$ is a 
dominant morphism.
\end{lemma}

\begin{proof}
Let $M \subset \Gamma(S, \mathcal{O}_S)$ be the submodule corresponding
to the largest quasi-coherent $\mathcal{O}_S$-module contained 
in the kernel of $f^\sharp$. Then any element $a \in M$ is mapped
to zero by $f^\sharp$. However, $f^\sharp(a)$ is the element of
$$
\Gamma(S, f_*\mathcal{O}_X) = \Gamma(X, \mathcal{O}_X) =
\Gamma(S, \mathcal{O}_S)
$$
corresponding to $a$ itself! Thus $a = 0$. Hence $M = 0$ which proves
the first assertion. Note that this is equivalent to the morphism 
$f : X \to S$ being scheme-theoretically surjective.

\medskip\noindent
If $X$ is quasi-compact, then $\Ker(f^\sharp)$ is quasi-coherent by
Morphisms, Lemma \ref{morphisms-lemma-quasi-compact-scheme-theoretic-image}.
Hence $\Ker(f^\sharp) = 0$ and $f^\sharp$ is injective.
In this case, $f$ is a dominant morphism by part (4) of
Morphisms, Lemma \ref{morphisms-lemma-quasi-compact-scheme-theoretic-image}.
\end{proof}

\end{document}
